% -------------------------------------------------------------------------
% WBS ID: RES-VER-SPEC
% Deliverable: Integrated Simulation Test Specification
% Status: In progress
% Evidence status: V1--V2--V3 hierarchy established; execution pending
% Review status: Not reviewed
% Last updated: 2026-07-19
% Dependencies: RES-MOD-SPEC, RES-NUM-SPEC, RES-DAT-CAS, RES-GEO-SPEC
% Downstream interfaces: Software test implementation, validation campaign
% -------------------------------------------------------------------------

\section{RES-VER-SPEC --- Integrated Simulation Test Specification}
\label{sec:res-ver-spec}

\subsection{Purpose and Scope}
\label{sec:res-ver-spec-purpose}

This Deliverable integrates the verification and validation plan for
the regional 2D NLSWE model, the local 3D URANS--VOF fixed-barrier
model and the regional-to-local replay interface. It defines required
tests and emitted metrics only; it does not claim that validation has
been executed.

\subsection{Research Questions}
\label{sec:res-ver-spec-research-questions}

% TODO: State the questions that must be answered before the Deliverable
% can be accepted.

\subsection{Terminology and Definitions}
\label{sec:res-ver-spec-definitions}

% TODO: Define the technical terminology, symbols, quantities and
% classifications used in this Deliverable.

\subsection{Literature Evidence}
\label{sec:res-ver-spec-literature-evidence}

\textcite{grilli2013tohoku} and \textcite{wei2013japan} support a
Tohoku hindcast stage in which source, propagation and coastal
hydrodynamic observations are compared before local impact claims are
made. \textcite{qin2018comparison} and \textcite{watanabe2022validation}
support local validation using velocity, pressure and force outputs
rather than free-surface elevation alone. Project conclusion: the
test specification shall be staged as V1 solver capability, V2 Tohoku
hydrodynamics and V3 local impact/loading.

\subsection{Governing Relations and Quantitative Definitions}
\label{sec:res-ver-spec-governing-relations}

% TODO: Record relevant equations, dimensionless groups, empirical models,
% extraction rules or quantitative definitions.
%
% Use align environments for displayed equations.
% Every displayed equation must have a unique \label{eq:...}.
% Do not place punctuation after displayed equations.

\subsection{Test Objective}
\label{sec:res-ver-spec-test-objective}

The integrated test objective is to verify that:

\begin{enumerate}
  \item V1 solver-capability tests pass for regional propagation,
    local free-surface flow, wall/obstacle impact and conservation;
  \item V2 Tohoku hydrodynamic tests compare source-to-coast
    propagation, arrival time, waveform, run-up and inundation against
    accepted observations for the Kamaishi and Sendai corridors;
  \item V3 impact/loading tests compare local pressure, force, moment,
    impulse, overtopping and energy outputs for fixed rigid barrier
    classes.
\end{enumerate}

The locked 2D NLSWE finite-volume specification must preserve
conservation, positivity, well-balanced still water, moving-bed source
behaviour, wet/dry shoreline motion, boundary non-reflection and
Tohoku-scale propagation before it can drive accepted replay histories.

\subsection{Reference or Benchmark Solution}
\label{sec:res-ver-spec-reference-benchmark-solution}

% TODO: Complete this RES-VER Domain-specific subsection.

\subsection{Test Configuration}
\label{sec:res-ver-spec-test-configuration}

% TODO: Complete this RES-VER Domain-specific subsection.

\subsection{Measured Quantities}
\label{sec:res-ver-spec-measured-quantities}

Regional measured quantities shall include:

\begin{itemize}
  \item water depth $h$ and free-surface elevation $\eta$;
  \item depth-integrated discharges $q_x$ and $q_y$;
  \item arrival time, peak elevation, phase and dominant period;
  \item run-up elevation, inundation distance and wet/dry extent;
  \item mass balance and momentum-balance diagnostics;
  \item boundary-reflection indicators and sponge-layer attenuation;
  \item mesh-convergence and CFL-sensitivity measures.
\end{itemize}

Local and interface measured quantities shall include:

\begin{itemize}
  \item $\alpha$ boundedness, phase volume and mass conservation;
  \item local CFL, pressure-correction residual and rejected timestep
    count;
  \item replay-boundary water level, velocity, discharge and pressure
    state;
  \item local pressure and shear histories on barrier patches;
  \item resultant force and moment histories;
  \item force impulse, pressure impulse and centre of pressure;
  \item run-up, overtopping and transmitted/reflected energy metrics.
\end{itemize}

\subsection{Error Metrics}
\label{sec:res-ver-spec-error-metrics}

% TODO: Complete this RES-VER Domain-specific subsection.

\subsection{Tolerance and Acceptance Criteria}
\label{sec:res-ver-spec-tolerance-acceptance-criteria}

% TODO: Complete this RES-VER Domain-specific subsection.

\subsection{Execution Handoff}
\label{sec:res-ver-spec-execution-handoff}

The following V1--V2--V3 verification and validation sequence is
required before validation is declared complete:

\begin{enumerate}
  \item V1 regional lake at rest over non-flat bathymetry;
  \item V1 one- and two-dimensional dam-break or bore propagation;
  \item V1 wet/dry run-up with positivity and mass-balance audit;
  \item V1 local VOF boundedness, pressure-correction residual and
    wall/obstacle impact benchmarks;
  \item V1 moving-bottom tsunami-generation test;
  \item V1 open/radiation boundary and sponge-reflection test;
  \item V1 mesh refinement and CFL sensitivity on regional and local
    domains;
  \item V2 Tohoku waveform/current/run-up or inundation comparison using
    data not already consumed by source calibration.
  \item V2 Kamaishi and Sendai replay-boundary consistency checks.
  \item V3 local fixed-barrier run-up, overtopping, pressure, force,
    moment, impulse and energy-metric tests.
\end{enumerate}

Execution evidence belongs to the verification and Software workstreams
and shall not be inferred merely from completion of the research
specification.

\subsection{Candidate Approaches}
\label{sec:res-ver-spec-candidate-approaches}

% TODO: Define the credible candidate approaches considered.

\subsection{Critical Comparison}
\label{sec:res-ver-spec-critical-comparison}

% TODO: Compare candidates using physically and project-relevant criteria.
% Do not select an approach solely on popularity or implementation ease.

\subsection{Assumptions and Applicability}
\label{sec:res-ver-spec-assumptions}

% TODO: State assumptions and the conditions under which they are valid.

\subsection{Limitations and Failure Conditions}
\label{sec:res-ver-spec-limitations}

% TODO: State where the evidence, model or selected approach ceases to be
% reliable or applicable.

\subsection{Project-Specific Conclusions}
\label{sec:res-ver-spec-conclusions}

% TODO: Convert the evidence into conclusions specific to the Studentship
% project. Do not merely restate literature findings.

The regional mathematical and numerical model is ready for verification
planning, and the local fixed-rigid barrier model now has a staged
acceptance path. Neither is yet a validated solver or validated Tohoku
case-study simulation.

\subsection{Selected Approach and Rationale}
\label{sec:res-ver-spec-selected-approach}

% TODO: Record the selected approach only after sufficient evidence exists.
% State the alternatives rejected and the reasons for rejection.

\subsection{Unresolved Decisions}
\label{sec:res-ver-spec-unresolved-decisions}

% TODO: Record unresolved questions, required evidence and decision owners.

Open verification decisions include final benchmark tolerances, final
case-study observation sets, calibration/validation split, final mesh
refinement levels, and acceptance thresholds for boundary reflection and
wet/dry mass-balance corrections. Local open decisions include final
$\alpha$ boundedness tolerance, CFL limit, pressure-residual
tolerance, force/moment consistency tolerance and energy-diagnostic
normalisation.

\subsection{Requirements and Handoffs}
\label{sec:res-ver-spec-requirements}

Software shall emit the RES-VER-MET metrics for every run and tag each
report as V1, V2 or V3. RES-DAT-CAS supplies observation bundles and
calibration/validation split metadata; RES-GEO-SPEC supplies geometry
identifiers; RES-NUM-SPEC supplies CFL, residual, boundedness and
mass-conservation diagnostics. This handoff supports test
implementation but does not itself execute validation.

\subsection{Deliverable Completion Record}
\label{sec:res-ver-spec-completion-record}

% TODO: Record whether the Deliverable is:
%
% - Not started
% - In progress
% - Draft complete
% - Reviewed
% - Accepted
%
% Acceptance requires:
%
% - the research questions to be answered;
% - claims to be supported by appropriate evidence;
% - assumptions and limitations to be explicit;
% - the project-specific conclusion to be stated;
% - downstream requirements to be traceable;
% - unresolved decisions to be closed or formally deferred.
